% GENERATED FILE. Edit canonical lessons, then run npm run generate:books.
% canonical-source-hash: fnv1a64:378d7c61f6355369
% canonical-lessons: JA-C370-matomaru, JA-C370-tagaini, JA-C370-onaji, JA-C370-daitai, JA-C370-pittari

\chapter{Reaching Agreement}
\label{ch:ja-reaching-agreement}
\hlchaptermodality{370}

\begin{chapteropening}
\textbf{By the end of this chapter:} \emph{I can say what we will do once we agree.}

Complete the last lesson of chapter 370: I can say what we will do once we agree.
\end{chapteropening}

\section[\texorpdfstring{\ja{まとまる} (matomaru)}{まとまる (matomaru)}]{\ja{まとまる} (matomaru) — to be settled; to come together}

\label{lesson:JA-C370-matomaru}



\begin{quote}
Before the new one: say the Japanese for a demand, then the Japanese for a policy.
\end{quote}

\subsection*{You'll want to know: \ja{まとまる}}
\textbf{\ja{まとまる}} — \emph{matomaru} — \textquotedblleft{}to be settled; to come together\textquotedblright{}.

\begin{quote}
\textbf{\ja{いけんが} \ja{まとまったら、ていあんして} \ja{ください。}} — \emph{iken ga matomattara, teian shite kudasai} — \textquotedblleft{}Once your ideas come together, please make a proposal.\textquotedblright{}
\end{quote}

Polite \textbf{\ja{まとまります}}, \emph{matomarimasu}; a godan verb. Past (polite): \textbf{\ja{いけんが} \ja{まとまりました。}}, \emph{iken ga matomarimashita}.

\begin{cousinweb}
\textbf{\ja{まとまる}} (\emph{matomaru}) is a native Japanese word (wago), inherited rather than borrowed from Chinese, and its older origin is not settled. JMdict notes that it is usually written in kana (JMdict).
\end{cousinweb}

\begin{grammarlens}[title={What you've built}]
One more word for saying what happens if.
\end{grammarlens}

\subsection*{Your turn}
\begin{itemize}
\raggedright
  \item \emph{Say it:} \emph{matomaru}
  \item \emph{Say it:} \emph{matomaru}, once more
  \item \emph{Say it:} \emph{iken ga matomattara, teian shite kudasai} — Once your ideas come together, please make a proposal.
  \item \emph{Recall:} say the Japanese for a demand, then the Japanese for a policy, then say \emph{matomaru} again
\end{itemize}

\begin{tcolorbox}[breakable,colback=teal!4,colframe=teal!35!black,title={Before you move on}]
What does \ja{まとまる} mean? (\textquotedblleft{}To be settled; to come together\textquotedblright{}.) Say it once more.
\end{tcolorbox}
\section[\texorpdfstring{\ja{たがいに} (tagaini)}{たがいに (tagaini)}]{\ja{たがいに} (tagaini) — each other; mutually}

\label{lesson:JA-C370-tagaini}



\begin{quote}
Before the new one: say the Japanese for a policy, then the Japanese for to be settled.
\end{quote}

\subsection*{You'll want to know: \ja{たがいに}}
\textbf{\ja{たがいに}} — \emph{tagaini} — \textquotedblleft{}each other; mutually\textquotedblright{}.

\begin{quote}
\textbf{\ja{たがいに} \ja{たすけたら、はやく} \ja{おわるでしょう。}} — \emph{tagaini tasuketara, hayaku owaru deshō} — \textquotedblleft{}If we help each other, we will probably finish sooner.\textquotedblright{}
\end{quote}

\begin{cousinweb}
\textbf{\ja{たがいに}} (\emph{tagaini}) is a native Japanese word (wago), inherited rather than borrowed from Chinese, and its older origin is not settled. Its kanji was attached to it later as a kun'yomi, a native reading given to a Chinese character (JMdict, KANJIDIC2).
\end{cousinweb}

\begin{grammarlens}[title={What you've built}]
One more word for saying what happens if.
\end{grammarlens}

\subsection*{Your turn}
\begin{itemize}
\raggedright
  \item \emph{Say it:} \emph{tagaini}
  \item \emph{Say it:} \emph{tagaini}, once more
  \item \emph{Say it:} \emph{tagaini tasuketara, hayaku owaru deshō} — If we help each other, we will probably finish sooner.
  \item \emph{Recall:} say the Japanese for a policy, then the Japanese for to be settled, then say \emph{tagaini} again
\end{itemize}

\begin{tcolorbox}[breakable,colback=teal!4,colframe=teal!35!black,title={Before you move on}]
What does \ja{たがいに} mean? (\textquotedblleft{}Each other; mutually\textquotedblright{}.) Say it once more.
\end{tcolorbox}
\section[\texorpdfstring{\ja{おなじ} (onaji)}{おなじ (onaji)}]{\ja{おなじ} (onaji) — the same}

\label{lesson:JA-C370-onaji}



\begin{quote}
Before the new one: say the Japanese for to be settled, then the Japanese for each other.
\end{quote}

\subsection*{You'll want to know: \ja{おなじ}}
\textbf{\ja{おなじ}} — \emph{onaji} — \textquotedblleft{}the same\textquotedblright{}.

\begin{quote}
\textbf{\ja{みんなが} \ja{おなじ} \ja{いけんに} \ja{なったら、きめます。}} — \emph{minna ga onaji iken ni nattara, kimemasu} — \textquotedblleft{}Once everyone has the same opinion, we will decide.\textquotedblright{}
\end{quote}

\begin{cousinweb}
\textbf{\ja{おなじ}} (\emph{onaji}) is a native Japanese word (wago), inherited rather than borrowed from Chinese, and its older origin is not settled. Its kanji was attached to it later as a kun'yomi, a native reading given to a Chinese character (JMdict, KANJIDIC2).
\end{cousinweb}

\begin{grammarlens}[title={What you've built}]
One more word for saying what happens if.
\end{grammarlens}

\subsection*{Your turn}
\begin{itemize}
\raggedright
  \item \emph{Say it:} \emph{onaji}
  \item \emph{Say it:} \emph{onaji}, once more
  \item \emph{Say it:} \emph{minna ga onaji iken ni nattara, kimemasu} — Once everyone has the same opinion, we will decide.
  \item \emph{Recall:} say the Japanese for to be settled, then the Japanese for each other, then say \emph{onaji} again
\end{itemize}

\begin{tcolorbox}[breakable,colback=teal!4,colframe=teal!35!black,title={Before you move on}]
What does \ja{おなじ} mean? (\textquotedblleft{}The same\textquotedblright{}.) Say it once more.
\end{tcolorbox}
\section[\texorpdfstring{\ja{だいたい} (daitai)}{だいたい (daitai)}]{\ja{だいたい} (daitai) — roughly; mostly}

\label{lesson:JA-C370-daitai}



\begin{quote}
Before the new one: say the Japanese for each other, then the Japanese for the same.
\end{quote}

\subsection*{You'll want to know: \ja{だいたい}}
\textbf{\ja{だいたい}} — \emph{daitai} — \textquotedblleft{}roughly; mostly\textquotedblright{}.

\begin{quote}
\textbf{\ja{だいたい} \ja{わかったら、はじめて} \ja{ください。}} — \emph{daitai wakattara, hajimete kudasai} — \textquotedblleft{}Once you more or less understand, please begin.\textquotedblright{}
\end{quote}

\begin{cousinweb}
\textbf{\ja{だいたい}} (\emph{daitai}) is Sino-Japanese (kango): two Chinese characters read with their on'yomi, readings borrowed from Chinese. They are \emph{dai}, \textquotedblleft{}large\textquotedblright{} (Middle Chinese \emph{dajH}) and \emph{tai}, \textquotedblleft{}body\textquotedblright{} (\emph{thejX}) (KANJIDIC2; Middle Chinese in Baxter's transcription).
\end{cousinweb}

\begin{grammarlens}[title={What you've built}]
One more word for saying what happens if.
\end{grammarlens}

\subsection*{Your turn}
\begin{itemize}
\raggedright
  \item \emph{Say it:} \emph{daitai}
  \item \emph{Say it:} \emph{daitai}, once more
  \item \emph{Say it:} \emph{daitai wakattara, hajimete kudasai} — Once you more or less understand, please begin.
  \item \emph{Recall:} say the Japanese for each other, then the Japanese for the same, then say \emph{daitai} again
\end{itemize}

\begin{tcolorbox}[breakable,colback=teal!4,colframe=teal!35!black,title={Before you move on}]
What does \ja{だいたい} mean? (\textquotedblleft{}Roughly; mostly\textquotedblright{}.) Say it once more.
\end{tcolorbox}
\section[\texorpdfstring{\ja{ぴったり} (pittari)}{ぴったり (pittari)}]{\ja{ぴったり} (pittari) — perfectly; exactly}

\label{lesson:JA-C370-pittari}



\begin{quote}
Before the new one: say the Japanese for the same, then the Japanese for roughly.
\end{quote}

\subsection*{You'll want to know: \ja{ぴったり}}
\textbf{\ja{ぴったり}} — \emph{pittari} — \textquotedblleft{}perfectly; exactly\textquotedblright{}.

\begin{quote}
\textbf{\ja{けいさんが} \ja{ぴったり} \ja{あったら、おわりです。}} — \emph{keisan ga pittari attara, owari desu} — \textquotedblleft{}If the sums come out exactly, we are done.\textquotedblright{}
\end{quote}

\begin{cousinweb}
\textbf{\ja{ぴったり}} (\emph{pittari}) is a native mimetic word, one of the Japanese words that picture a manner or a feeling in sound; JMdict marks it as mimetic. It is inherited, not borrowed, and its older origin is not settled (JMdict).
\end{cousinweb}

\begin{grammarlens}[title={What you've built}]
One more word for saying what happens if.
\end{grammarlens}

\subsection*{Your turn}
\begin{itemize}
\raggedright
  \item \emph{Say it:} \emph{pittari}
  \item \emph{Say it:} \emph{pittari}, once more
  \item \emph{Say it:} \emph{keisan ga pittari attara, owari desu} — If the sums come out exactly, we are done.
  \item \emph{Recall:} say the Japanese for the same, then the Japanese for roughly, then say \emph{pittari} again
\end{itemize}

\begin{tcolorbox}[breakable,colback=teal!4,colframe=teal!35!black,title={Before you move on}]
What does \ja{ぴったり} mean? (\textquotedblleft{}Perfectly; exactly\textquotedblright{}.) Say it once more.
\end{tcolorbox}
